% =============================================================
% Capítulo 2 de la tesis:
% Construcción de la intersección de una PRCG con un lenguaje regular
% =============================================================
% Pensado para compilarse de manera independiente con pdflatex.
% Para integrarlo en la tesis completa, basta con copiar el contenido
% que va entre \begin{document} y \end{document} y, en la tesis
% principal, asegurarse de cargar los paquetes que aparecen en el
% preámbulo.
% =============================================================

\documentclass[11pt, letterpaper]{article}

% --- Codificación e idioma ---
\usepackage[utf8]{inputenc}
\usepackage[T1]{fontenc}
\usepackage[spanish, provide=*]{babel}

% --- Matemáticas ---
\usepackage{amsmath}
\usepackage{amssymb}
\usepackage{amsthm}
\usepackage{mathtools}

% --- Layout ---
\usepackage[margin=1in]{geometry}
\usepackage{parskip}
\setlength{\parindent}{1.5em}
\setlength{\parskip}{0.5ex plus 0.2ex minus 0.1ex}

% --- Tipografía y enlaces ---
\usepackage{microtype}
\usepackage[hidelinks]{hyperref}

% --- Bloques de "código" (gramáticas y autómatas) ---
\usepackage{fancyvrb}
\usepackage{tcolorbox}

% --- Entornos para definiciones, teoremas, observaciones, condiciones ---
\theoremstyle{definition}
\newtheorem{definicion}{Definición}[section]
\newtheorem{condicion}{Condición}[section]
\theoremstyle{plain}
\newtheorem{teorema}{Teorema}[section]
\newtheorem{observacion}{Observación}[section]
\newtheorem{afirmacion}{Afirmación}[section]

% Renombrar las palabras que babel pone en negrita
\renewcommand{\proofname}{Demostración}

% --- Citas pendientes: que [REF] aparezca textualmente ---
% Cuando llegue el momento de añadir la bibliografía real, cambiar
% \cite por su definición estándar y añadir \bibliography{...} al final.
\renewcommand{\cite}[1]{[\textsc{#1}]}

% --- Macros del capítulo ---
\newcommand{\PRCG}{\mbox{PRCG}}
\newcommand{\sRCG}{\mbox{sRCG}}
\newcommand{\AFD}{\mbox{AFD}}
\newcommand{\states}{\mathcal{Q}}      % Q (estados)
\newcommand{\statesT}{Q}                % Q en uso normal
\newcommand{\final}{F}
\newcommand{\initial}{q_0}
\newcommand{\Lang}{L}                   % L(G), L(A)
\newcommand{\Tnt}{T}                    % alfabeto de terminales
\newcommand{\Vvar}{V}                   % conjunto de variables
\newcommand{\Pset}{P}                   % conjunto de cláusulas
\newcommand{\Nset}{N}                   % no terminales
\newcommand{\trans}{\delta}
\newcommand{\transit}{\delta^{*}}
\newcommand{\eps}{\varepsilon}
\newcommand{\concat}[2]{#1 #2}
\DeclareMathOperator{\arity}{ar}

% =============================================================
\begin{document}

% =============================================================
\section*{Capítulo 2 \\[0.3em] Construcción de la intersección de una PRCG con un lenguaje regular}
\addcontentsline{toc}{section}{Capítulo 2}
\setcounter{section}{2}
\setcounter{subsection}{0}

% =============================================================
\subsection{Introducción del capítulo}

En el capítulo anterior se presentó el formalismo de las gramáticas de concatenación de rango y se dejó planteado el problema central de la tesis: dada una \PRCG\ $G$ y un autómata finito $A$ que reconoce un lenguaje regular $R$, construir explícitamente un objeto que represente $\Lang(G) \cap R$.

Este capítulo aporta dicha construcción. La salida es una nueva \PRCG, denotada $G'$, cuyo lenguaje es exactamente $\Lang(G) \cap R$. La construcción se realiza por composición sintáctica de las cláusulas de $G$ con la estructura de transiciones de $A$, y produce un objeto cuya cantidad de símbolos no terminales puede crecer exponencialmente en el tamaño de $A$ y en la aridad de los predicados de $G$. Esa cota exponencial no se debe a una elección de implementación, sino a la cota inferior establecida por Bertsch y Nederhof~\cite{REF}, según la cual ningún algoritmo polinomial general puede existir para esta tarea, salvo que $\mathbf{P} = \mathbf{NP}$.

La construcción se inspira en la prueba del Teorema 3.9 de Seki et al.~\cite{REF} para la intersección de gramáticas multi-contextuales con lenguajes regulares, y en la construcción polinomial que Bertsch y Nederhof~\cite{REF} dan para la subclase de las \sRCG. Ninguna de las dos referencias trata el caso general de las \PRCG, donde la no-linealidad permite copia y borrado entre los argumentos de una cláusula. Esa diferencia es la que obliga a explicitar una condición de coherencia que en la construcción polinomial para \sRCG\ está implícita en la forma en que se parametriza la gramática producto, pero que en \PRCG\ general no puede dejarse implícita: cuando una misma variable aparece en dos o más posiciones de una cláusula, las instancias que se generan deben asignarle el mismo par de estados del autómata. Esa condición se nombra \textbf{C2} y se formaliza en la sección~\ref{sec:construccion}. La verificación operacional reportada en la sección~\ref{sec:literatura} muestra que la condición es necesaria incluso en gramáticas lineales en el sentido de Boullier, y que su rol pleno se manifiesta cuando una variable se repite en cabeza y cuerpo de manera no trivial, característica propia de las \PRCG\ con copia genuina.

El capítulo se organiza en seis secciones. En la sección~\ref{sec:preliminares} se recuerdan las definiciones de \PRCG, autómata finito y rango sobre un autómata. En la sección~\ref{sec:idea} se presenta la idea informal de la construcción mediante un ejemplo que muestra cómo los rangos clásicos sobre la cadena de entrada se reemplazan por pares de estados del autómata. En la sección~\ref{sec:construccion} se da la construcción formal de $G'$ y se demuestra el teorema de corrección. En la sección~\ref{sec:tamano} se analiza el tamaño de $G'$ y se muestra un ejemplo trabajado completo. En la sección~\ref{sec:literatura} se sitúa la construcción frente a la literatura, se explica el papel que juega la semántica de \PRCG\ en la suficiencia de la condición \textbf{C2} y se reporta una verificación operacional. En la sección~\ref{sec:transicion} se discute brevemente la relación de $G'$ con el problema del vacío, que se aborda en el capítulo siguiente.

% =============================================================
\subsection{Preliminares}\label{sec:preliminares}

Esta sección fija la notación. Los conceptos ya introducidos en el capítulo~1 se reproducen aquí en la forma exacta que se utiliza en las pruebas de las secciones siguientes.

\subsubsection{Gramáticas de concatenación de rango positivas}

Una \emph{gramática de concatenación de rango positiva}, abreviada \PRCG, es una tupla $G = (\Nset, \Tnt, \Vvar, \Pset, S)$, donde $\Nset$ es un conjunto finito de predicados, cada uno con una aridad fija; $\Tnt$ es un alfabeto finito de símbolos terminales; $\Vvar$ es un conjunto finito de variables; $\Pset$ es un conjunto finito de cláusulas; y $S \in \Nset$ es el predicado distinguido, de aridad uno.

Una cláusula de $\Pset$ tiene la forma
\[
A_0(\alpha_{0,1}, \dots, \alpha_{0, a_0}) \to A_1(\alpha_{1,1}, \dots, \alpha_{1, a_1}) \; \cdots \; A_m(\alpha_{m,1}, \dots, \alpha_{m, a_m}),
\]
donde cada $A_j \in \Nset$ tiene aridad $a_j$, y cada $\alpha_{j,i} \in (\Tnt \cup \Vvar)^{*}$ es una sucesión de símbolos terminales y variables. La parte izquierda de la flecha se denomina cabeza y la parte derecha se denomina cuerpo. Cuando $m = 0$, la cláusula se llama terminal y se escribe simplemente $A_0(\alpha_{0,1}, \dots, \alpha_{0, a_0}) \to \eps$.

La semántica de las \PRCG\ opera por sustitución de rangos. Un \emph{rango} sobre una cadena $w \in \Tnt^{*}$ es un par $\langle i, j \rangle$ con $0 \leq i \leq j \leq \lvert w \rvert$. El rango $\langle i, j \rangle$ representa la subcadena $w[i+1 \ldots j]$. Una \emph{instancia de cláusula} sobre $w$ es la cláusula que resulta de reemplazar cada variable que aparece en $\alpha_{j,i}$ por un rango sobre $w$, de modo coherente: dos ocurrencias de la misma variable reciben el mismo rango. Una instancia es válida si la concatenación de los rangos asignados a $\alpha_{j,i}$ coincide con un rango único, y si los símbolos terminales que aparecen en $\alpha_{j,i}$ corresponden a las posiciones de $w$ que el rango cubre. El lenguaje generado por $G$, denotado $\Lang(G)$, es el conjunto de cadenas $w \in \Tnt^{*}$ tales que $S(\langle 0, \lvert w \rvert \rangle)$ puede derivarse a la cadena vacía mediante la aplicación de instancias válidas de cláusulas.

\subsubsection{Hipótesis sobre la clase de \PRCG\ tratada}\label{sec:hipotesis}

El algoritmo del presente capítulo trata las \PRCG\ \emph{top-down non-erasing}, en el sentido de Boullier~\cite{REF}, Definición~8, página~11: una \PRCG\ $G$ es top-down non-erasing si, para toda cláusula $c \in \Pset$, cada variable que aparece en la cabeza de $c$ aparece también en el cuerpo de $c$. La restricción es puramente sintáctica y se verifica por inspección directa de las cláusulas.

Esta clase es estrictamente más amplia que las \sRCG. Una \sRCG, según la Definición~9 de Boullier~\cite{REF}, página~13, exige tres propiedades simultáneas: linealidad, non-erasure en ambas direcciones y no-combinatorialidad de los argumentos del cuerpo. La hipótesis top-down non-erasing exige solamente la segunda en su dirección descendente. En particular, una \PRCG\ top-down non-erasing puede contener cláusulas no-lineales, en las que una variable aparece más de una vez en cabeza o cuerpo, y cláusulas combinatorias, en las que un argumento del cuerpo contiene terminales mezclados con variables.

La restricción no compromete la generalidad del algoritmo. Por la cadena de Properties~1, 2 y 3 demostradas por Boullier~\cite{REF}, páginas~11--12, toda \PRCG\ admite una \PRCG\ equivalente que es no-combinatorial y non-erasing. La cadena se compone de tres transformaciones encadenadas sobre las cláusulas: la Property~1 elimina la combinatorialidad de los argumentos del cuerpo, la Property~2 elimina el borrado ascendente, y la Property~3 elimina el borrado descendente mediante la introducción de un predicado auxiliar $\mathrm{Any}$ que reconoce $\Tnt^{*}$. La gramática equivalente que resulta de la cadena es no-combinatorial y non-erasing en ambos sentidos, y en particular top-down non-erasing, lo que la hace tratable por el algoritmo del presente capítulo. La transformación queda fuera del alcance del capítulo y se invoca solo como precondición externa cuando la entrada original no cumple la hipótesis.

Sobre la linealidad, Boullier~\cite{REF}, página~12, demuestra explícitamente que ninguna gramática lineal es equivalente, en general, a una no-lineal: \emph{``these equivalence properties do not hold for linearity (i.e. in the general case, there is no linear grammar equivalent to a non-linear one). In other words, non-linearity brings some formal power to the formalism.''} Las \PRCG\ no-lineales son por tanto un objeto estrictamente más expresivo que las \sRCG, y la diferencia es la fuente de la cerradura de las \PRCG\ bajo intersección y de las propiedades que se demuestran en este capítulo. El algoritmo aporta una construcción explícita de la \PRCG\ producto en una clase de \PRCG\ que cubre toda \PRCG\ hasta equivalencia de lenguajes, y que es estrictamente más expresiva que la clase \sRCG\ tratada por Bertsch y Nederhof~\cite{REF}.

\subsubsection{Autómatas finitos deterministas}

Un \emph{autómata finito determinista}, abreviado \AFD, es una tupla $A = (\statesT, \Sigma, \trans, \initial, \final)$, donde $\statesT$ es un conjunto finito de estados; $\Sigma$ es un alfabeto, que se asume igual al alfabeto $\Tnt$ de la \PRCG; $\trans : \statesT \times \Sigma \to \statesT$ es la función de transición; $\initial \in \statesT$ es el estado inicial; y $\final \subseteq \statesT$ es el conjunto de estados finales. La función de transición se extiende a $\transit : \statesT \times \Sigma^{*} \to \statesT$ en la forma habitual. El lenguaje reconocido por $A$, denotado $\Lang(A)$, es el conjunto de cadenas $w \in \Sigma^{*}$ tales que $\transit(\initial, w) \in \final$.

En todo el capítulo se asume que el autómata $A$ está fijado de antemano y que sus componentes $\statesT$, $\trans$, $\initial$ y $\final$ son entrada de la construcción. La definición de $\trans$ excluye las transiciones $\eps$: la función está tipada sobre $\statesT \times \Sigma$ y no sobre $\statesT \times (\Sigma \cup \{\eps\})$. La restricción es importante para la corrección de la construcción Bar-Hillel-style. Pasti, Opedal, Pimentel, Eisner, Cotterell y Vieira~\cite{REF} demostraron que la construcción Bar-Hillel para gramáticas libres del contexto deja de funcionar cuando el autómata admite transiciones $\eps$, y propusieron una corrección específica para ese caso. La construcción del presente capítulo no aborda la extensión a autómatas con transiciones $\eps$; un autómata regular cualquiera con transiciones $\eps$ puede convertirse a la forma sin transiciones $\eps$ mediante el procedimiento clásico, sin pérdida de generalidad.

\subsubsection{Rangos sobre un autómata}

La idea central de la construcción consiste en sustituir el papel que tienen los rangos clásicos por una versión que se ajuste a un autómata. Un \emph{rango-autómata} es un par de estados $\langle q, q' \rangle \in \statesT \times \statesT$. La interpretación informal de $\langle q, q' \rangle$ es la siguiente: existe una cadena $u \in \Sigma^{*}$ tal que $\transit(q, u) = q'$. La cadena $u$ no se especifica; queda implícita en el par.

Cuando $A$ es el autómata canónico que reconoce todas las cadenas de longitud arbitraria sobre $\Sigma$, los rangos-autómata se corresponden con rangos clásicos. En el caso general, el conjunto de cadenas que conectan dos estados puede tener tamaño infinito o ser vacío, lo que constituye precisamente la diferencia entre un autómata cualquiera y la cadena fija sobre la que opera la semántica clásica de las \PRCG.

Un \emph{vector de rangos-autómata} de longitud $k$ es una tupla $\vec{q} = (\langle q_1, q'_1 \rangle, \dots, \langle q_k, q'_k \rangle) \in (\statesT \times \statesT)^{k}$. Estos vectores se utilizan para indexar los predicados en la gramática producto que se construye en la sección siguiente.

% =============================================================
\subsection{Idea informal de la construcción}\label{sec:idea}

La construcción que se va a presentar parte de la siguiente observación: cuando una \PRCG\ opera sobre una cadena fija $w$, cada ocurrencia de cada variable se asocia a un rango $\langle i, j \rangle$ que apunta a una subcadena de $w$. Si en lugar de $w$ se considera el conjunto $R = \Lang(A)$ de todas las cadenas reconocidas por un autómata, ya no hay un único rango por variable, sino, potencialmente, varios pares de estados $\langle q, q' \rangle$ candidatos a representar las cadenas que cumplen los requisitos de la cláusula. La gramática producto debe enumerar todos esos candidatos y conservar solo los que respeten dos restricciones: la coherencia con las transiciones del autómata, y la igualdad entre las distintas ocurrencias de una misma variable.

El primer tipo de coherencia se hereda de la construcción de Bar-Hillel para gramáticas libres del contexto y de su extensión a \sRCG\ por Bertsch y Nederhof~\cite{REF}. Para cada cláusula de $G$ y cada elección de pares de estados que asigne uno a cada argumento de cada predicado, la cláusula producto solo se conserva si los pares se concatenan correctamente: si un argumento se descompone como $\alpha = \beta_1 \beta_2 \cdots \beta_l$, los estados intermedios entre los $\beta_j$ deben coincidir.

El segundo tipo de coherencia es la pieza estructural que la construcción Bar-Hillel-style necesita siempre que una variable aparece en más de una posición de una cláusula. Cuando una variable $X$ aparece dos o más veces en una cláusula, sea en la cabeza, en el cuerpo, o repartidas entre ambas partes, las ocurrencias deben asignarse al mismo par de estados. Esa restricción es trivial cuando la cláusula es enteramente lineal en cabeza y cuerpo en simultáneo, ya que la asignación natural por variable la cumple por construcción, pero la verificación operacional reportada en la sección~\ref{sec:literatura} muestra que omitirla, asignando un par independiente a cada ocurrencia sintáctica, produce sobre-generación incluso en \sRCG~\cite{REF}, esto es, lenguajes generados por la gramática producto que no pertenecen a $\Lang(G) \cap R$. La diferencia entre las \PRCG\ top-down non-erasing del presente capítulo y las \sRCG\ no está en la existencia de esta condición, sino en la diversidad de configuraciones donde actúa de manera no trivial: en \sRCG\ una variable puede aparecer a lo sumo una vez en la cabeza y una vez en el cuerpo, y la condición acopla únicamente esas dos ocurrencias; en una \PRCG\ top-down non-erasing una variable puede aparecer cualquier número de veces, y la condición acopla todas las ocurrencias simultáneamente.

Esa condición de coherencia tiene un alcance distinto al que tiene la construcción de Seki~\cite{REF} para gramáticas multi-contextuales paralelas, donde la copia se manifiesta en la función de composición de las reglas. En las \PRCG\ la copia se manifiesta en la repetición de variables dentro de una cláusula. La construcción que se presenta en la sección siguiente trata directamente esa repetición a nivel de las cláusulas producto.

\subsubsection{Ejemplo motivador}

Para fijar ideas, se considera la \PRCG\ $G_0$ con predicados $S$ de aridad uno y $\mathrm{Par}$ de aridad uno, sobre el alfabeto $\Tnt = \{a\}$, y las cláusulas siguientes:

\begin{Verbatim}[xleftmargin=2em]
  c1 :  S(X)        -> Par(X) Par(X)
  c2 :  Par(aaX)    -> Par(X)
  c3 :  Par(eps)    -> eps
\end{Verbatim}

El lenguaje $\Lang(G_0)$ es $\{ a^{2n} : n \geq 0 \}$. La cláusula $c_1$ contiene una variable $X$ que aparece dos veces dentro del cuerpo, una en cada uno de los predicados $\mathrm{Par}(X)$ del lado derecho. Eso constituye no-linealidad top-down en el sentido de Boullier~\cite{REF}: una variable aparece más de una vez en la unión de las tuplas de argumentos de los predicados del cuerpo. La no-linealidad de $c_1$ exige que las dos invocaciones de $\mathrm{Par}$ se hagan sobre la misma cadena, lo cual obliga a que $\mathrm{Par}(X)$ derive simultáneamente la cadena entera. Como $\mathrm{Par}$ por sí solo reconoce $\{a^{2n}\}$, la cláusula $c_1$ es semánticamente redundante respecto a $\mathrm{Par}$ aplicada una sola vez, pero la doble invocación es exactamente la situación donde la condición \textbf{C2} de la construcción que se presenta en este capítulo actúa de manera no trivial. Esto se hace explícito en la sección~\ref{sec:literatura}.

Sea ahora $A_0$ el \AFD\ con $\statesT = \{q_0, q_1, q_2\}$, alfabeto $\{a\}$, $q_0$ inicial, $\final = \{q_0\}$, y transiciones $\trans(q_0, a) = q_1$, $\trans(q_1, a) = q_2$, $\trans(q_2, a) = q_0$. El lenguaje $\Lang(A_0)$ es el conjunto de cadenas de $a$ con longitud múltiplo de $3$, esto es, $\{ a^{3k} : k \geq 0 \}$.

La intersección $\Lang(G_0) \cap \Lang(A_0)$ es el conjunto de cadenas de $a$ con longitud que sea simultáneamente múltiplo de $2$ y de $3$, esto es, múltiplo de $6$:
\[
\Lang(G_0) \cap \Lang(A_0) = \{ a^{6m} : m \geq 0 \}.
\]
La construcción que se desarrolla a continuación produce una \PRCG\ cuyo lenguaje es ese conjunto.

% =============================================================
\subsection{Construcción formal}\label{sec:construccion}

Esta sección define la \PRCG\ producto $G'$ y demuestra que su lenguaje es $\Lang(G) \cap \Lang(A)$.

\subsubsection{Definición de la gramática producto}

\begin{definicion}[Gramática producto $G'$]
\label{def:gprime}
Sean $G = (\Nset, \Tnt, \Vvar, \Pset, S)$ una \PRCG\ top-down non-erasing en el sentido de la sección~\ref{sec:hipotesis}, y $A = (\statesT, \Tnt, \trans, \initial, \final)$ un \AFD. La \emph{gramática producto} es la \PRCG
\[
G' = (\Nset', \Tnt, \Vvar, \Pset', S'),
\]
cuyos componentes se definen como sigue.
\end{definicion}

\paragraph{El conjunto de no terminales.} $\Nset'$ contiene un nuevo símbolo distinguido $S'$ de aridad uno y, por cada predicado $A \in \Nset$ de aridad $a$ y por cada vector $\vec{q} = (\langle q_1, q'_1 \rangle, \dots, \langle q_a, q'_a \rangle) \in (\statesT \times \statesT)^{a}$, un símbolo nuevo denotado $A[\vec{q}]$, con la misma aridad $a$ que $A$. Formalmente,
\[
\Nset' = \{S'\} \;\cup\; \{ A[\vec{q}] : A \in \Nset \text{ de aridad } a, \; \vec{q} \in (\statesT \times \statesT)^{a} \}.
\]

\paragraph{Las cláusulas iniciales.} Por cada estado final $q \in \final$, se añade una cláusula
\[
S'(X) \to S[\langle \initial, q \rangle](X).
\]
Estas cláusulas conectan el símbolo distinguido $S'$ con todas las maneras posibles de leer la cadena completa desde el estado inicial $\initial$ hasta un estado final.

\paragraph{Las cláusulas heredadas.} Por cada cláusula $c \in \Pset$ de $G$ con la forma
\[
A_0(\alpha_{0,1}, \dots, \alpha_{0, a_0}) \to A_1(\alpha_{1,1}, \dots, \alpha_{1, a_1}) \; \cdots \; A_m(\alpha_{m,1}, \dots, \alpha_{m, a_m}),
\]
y por cada asignación de pares de estados que satisfaga las dos condiciones que se enuncian a continuación, se añade a $\Pset'$ una cláusula heredada que se obtiene a partir de $c$ indexando cada predicado con el vector de pares de estados correspondiente.

La asignación que produce una cláusula heredada se especifica como sigue. Para cada predicado $A_j$ que aparece en $c$, sea $\vec{q}_j = (\langle p_{j,1}, p'_{j,1} \rangle, \dots, \langle p_{j, a_j}, p'_{j, a_j} \rangle)$ un vector de pares de estados de longitud $a_j$. La asignación se denota $\Phi$ y consiste en la familia $(\vec{q}_0, \vec{q}_1, \dots, \vec{q}_m)$.

\begin{condicion}[Coherencia interna de los argumentos]
\label{cond:c1}
Para cada argumento $\alpha_{j,i}$ que aparece en $c$, asignado al par $\langle p_{j,i}, p'_{j,i} \rangle$, sea $\alpha_{j,i} = \beta_1 \beta_2 \cdots \beta_l$, donde cada $\beta_u$ es o un símbolo terminal $a \in \Tnt$, o una variable $X \in \Vvar$. La condición \textbf{C1} establece que existe una sucesión de estados $r_0 = p_{j,i}, r_1, r_2, \dots, r_l = p'_{j,i}$ tal que: si $\beta_u = a \in \Tnt$, entonces $\trans(r_{u-1}, a) = r_u$; y si $\beta_u = X$ es una variable, entonces $r_{u-1}$ y $r_u$ coinciden con los estados que la asignación $\Phi$ ha asociado a la ocurrencia concreta de $X$ en $\alpha_{j,i}$.
\end{condicion}

Una vez fijada la asignación $\Phi$ a cada ocurrencia de cada variable, los estados intermedios entre los símbolos terminales y las variables quedan determinados por las transiciones de $A$. La condición \textbf{C1} se reduce, en la práctica, a verificar que las transiciones inducidas por los terminales son consistentes con los pares asignados.

\begin{condicion}[Coherencia entre ocurrencias de la misma variable]
\label{cond:c2}
Si una variable $X \in \Vvar$ aparece en dos o más posiciones de la cláusula $c$ ---sea en la cabeza, en el cuerpo, o repartidas entre ambas partes--- entonces los pares de estados que la condición \textbf{C1} asocia a todas esas ocurrencias deben ser idénticos entre sí.
\end{condicion}

La condición \textbf{C2} expresa que, durante una derivación de $G$, una variable $X$ se instancia en una única subcadena $u$, y esa única $u$ debe leerse en $A$ desde y hasta los mismos estados sin importar dónde aparezca $X$ en la cláusula. Sin la condición \textbf{C2} la gramática producto generaría cadenas en las que las ocurrencias de $X$ tomarían valores distintos, en contradicción con la semántica de $G$.

Cuando ambas condiciones \textbf{C1} y \textbf{C2} se cumplen para una asignación $\Phi$ aplicada a una cláusula $c$ de $G$, se añade a $\Pset'$ la cláusula heredada
\[
A_0[\vec{q}_0](\alpha_{0,1}, \dots, \alpha_{0, a_0}) \to A_1[\vec{q}_1](\alpha_{1,1}, \dots) \; \cdots \; A_m[\vec{q}_m](\alpha_{m,1}, \dots).
\]
Esa cláusula tiene la misma estructura sintáctica de $c$, con las mismas variables en las mismas posiciones, y difiere de $c$ únicamente en que cada predicado lleva ahora un vector de pares de estados como índice.

\paragraph{Cláusulas terminales.} Una cláusula terminal de $G$ tiene la forma $A(\alpha_{0,1}, \dots, \alpha_{0, a_0}) \to \eps$, donde cada $\alpha_{j,i}$ es una sucesión que solo contiene símbolos terminales o $\eps$. La ausencia de variables en los argumentos de una cláusula terminal es consecuencia directa de la hipótesis top-down non-erasing de la sección~\ref{sec:hipotesis}: una cláusula con cuerpo vacío no puede contener variables en la cabeza, ya que la hipótesis exige que toda variable de la cabeza aparezca en el cuerpo. Para cada vector $\vec{q}_0$ que satisfaga la condición \textbf{C1} sobre los argumentos terminales, se añade a $\Pset'$ la cláusula
\[
A[\vec{q}_0](\alpha_{0,1}, \dots, \alpha_{0, a_0}) \to \eps.
\]
En el caso particular en que algún argumento $\alpha_{j,i}$ es $\eps$, la condición \textbf{C1} obliga a que el par $\langle p_{j,i}, p'_{j,i} \rangle$ asociado satisfaga $p_{j,i} = p'_{j,i}$, ya que la cadena vacía no produce transición.

\subsubsection{Teorema de corrección}

\begin{teorema}[Corrección de la construcción]
\label{thm:correctness}
Sean $G$ una \PRCG\ top-down non-erasing y $A$ un \AFD\ sobre el mismo alfabeto $\Tnt$. Sea $G'$ la gramática producto de la Definición~\ref{def:gprime}. Entonces
\[
\Lang(G') = \Lang(G) \cap \Lang(A).
\]
\end{teorema}

\begin{proof}
La demostración se realiza por doble inclusión. La hipótesis de que $G$ es top-down non-erasing (sección~\ref{sec:hipotesis}) interviene en dos pasos específicos de la dirección $\Lang(G') \subseteq \Lang(G) \cap \Lang(A)$: en el caso base, donde garantiza que las cláusulas terminales no contienen variables, y en el paso inductivo, donde garantiza que las subcadenas asignadas a las variables de la cabeza quedan determinadas por las subderivaciones del cuerpo. Antes de proceder, se establece una observación auxiliar sobre el recorrido del autómata.

\begin{observacion}[Recorrido único del autómata sobre una cadena fija]
\label{obs:recorrido}
Sea $w \in \Tnt^{*}$ y $A$ un \AFD. La sucesión de estados $s_0, s_1, \dots, s_n$ con $s_0 = \initial$ y $s_k = \transit(\initial, w[1 \ldots k])$ está unívocamente determinada por $w$ y $A$, ya que $\trans$ es una función. Para cualquier rango $\langle i, j \rangle$ sobre $w$, el par $\langle s_i, s_j \rangle$ es el único par de estados consistente con el recorrido del autómata sobre la subcadena $w[i+1 \ldots j]$, y satisface $\transit(s_i, w[i+1 \ldots j]) = s_j$.
\end{observacion}

Esta observación se utiliza en ambas direcciones de la doble inclusión para asociar a cada rango sobre $w$ un par de estados unívocamente determinado.

\medskip
\noindent\textbf{Inclusión $\Lang(G') \subseteq \Lang(G) \cap \Lang(A)$.}\\
Sea $w \in \Lang(G')$. Por definición existe una derivación de $S'(\langle 0, \lvert w \rvert \rangle)$ hacia $\eps$ en $G'$. La primera regla aplicada es necesariamente una de las cláusulas iniciales $S'(X) \to S[\langle \initial, q \rangle](X)$, con $q \in \final$. Esa primera regla obliga a que la subderivación restante reconozca $w$ a partir del símbolo $S[\langle \initial, q \rangle]$ y, por unicidad del recorrido del autómata, fuerza $\transit(\initial, w) = q$, de donde $w \in \Lang(A)$.

Para concluir $w \in \Lang(G)$, se demuestra la siguiente afirmación inductiva, que es la versión reforzada de la propiedad estándar de corrección:

\begin{afirmacion}[Propiedad inductiva reforzada]
\label{afirm:induccion}
Si $A[\vec{q}]$ deriva en $G'$ una tupla de subcadenas $(u_1, \dots, u_a)$, entonces se cumplen las dos condiciones siguientes:
\begin{enumerate}
\item[(i)] $A$ deriva en $G$ la misma tupla $(u_1, \dots, u_a)$;
\item[(ii)] para cada componente $u_i$ de la tupla, $\transit(p_i, u_i) = p'_i$, donde $\langle p_i, p'_i \rangle$ es la $i$-ésima componente de $\vec{q}$.
\end{enumerate}
\end{afirmacion}

La afirmación se demuestra por inducción sobre la altura de la derivación de $A[\vec{q}]$ en $G'$.

\textit{Caso base.} La derivación tiene altura uno y procede aplicando una cláusula terminal $A[\vec{q}](\alpha_{0,1}, \dots, \alpha_{0, a}) \to \eps$. Por construcción de $G'$ esa cláusula está en $\Pset'$ porque la cláusula original $A(\alpha_{0,1}, \dots, \alpha_{0, a}) \to \eps$ está en $\Pset$ y la condición \textbf{C1} se satisface para $\vec{q}$. Por la hipótesis de que $G$ es top-down non-erasing, una cláusula con cuerpo vacío no contiene variables en la cabeza, ya que toda variable de la cabeza debe aparecer en el cuerpo. Por tanto cada $\alpha_{0,i}$ contiene únicamente símbolos terminales y, eventualmente, $\eps$. Las cadenas $u_i$ resultan de la concatenación de esos terminales. La condición (i) se cumple porque la cláusula original existe en $\Pset$ y aplicarla en $G$ produce la misma tupla. La condición (ii) se cumple porque la condición \textbf{C1}, aplicada en construcción, exige una sucesión de estados que va de $p_i$ a $p'_i$ siguiendo las transiciones de $A$ sobre los terminales de $\alpha_{0,i}$; esa sucesión es exactamente la lectura $\transit(p_i, u_i) = p'_i$.

\textit{Paso inductivo.} La derivación inicia con una cláusula heredada
\[
c' = A[\vec{q}](\alpha_{0,1}, \dots) \to A_1[\vec{q}_1](\alpha_{1,1}, \dots) \cdots A_m[\vec{q}_m](\alpha_{m,1}, \dots).
\]
Por construcción de $G'$ la cláusula original sin índices, $c$, está en $\Pset$, y la asignación $\Phi$ que asocia un par a cada variable que aparece en $c$ satisface las condiciones \textbf{C1} y \textbf{C2}.

Cada subderivación de $A_j[\vec{q}_j]$ tiene altura estrictamente menor. Por hipótesis de inducción, $A_j$ deriva en $G$ la tupla de subcadenas que la subderivación produce, y se satisface la propiedad de pares para cada componente.

Para construir la derivación correspondiente en $G$, se requiere asignar rangos a las variables que aparecen en $c$. Aquí interviene la hipótesis top-down non-erasing: toda variable que aparece en la cabeza de $c$ aparece también en el cuerpo de $c$, y por tanto las subderivaciones del cuerpo determinan los rangos asignados a las variables. Concretamente, una variable $X$ que aparece en el cuerpo lo hace dentro del argumento de algún predicado $A_j$; la subderivación de $A_j$ produce una tupla de subcadenas, y la subcadena correspondiente a la posición de $X$ es la asignada a $X$. La condición \textbf{C1} fija los pares de estados de los argumentos, lo que determina los puntos de corte entre terminales y variables. La cadena $u_i$ correspondiente a la $i$-ésima componente de la cabeza se obtiene concatenando, según el patrón $\alpha_{0,i}$, los terminales y las cadenas asignadas a las variables. Sin la hipótesis top-down non-erasing, una variable de la cabeza podría no aparecer en el cuerpo, y su rango quedaría indeterminado, lo cual impediría la construcción de la derivación correspondiente en $G$.

En este punto interviene la condición \textbf{C2} con su justificación semántica. Si una variable $X$ aparece en dos posiciones distintas de $c$, las dos ocurrencias en $c'$ tienen, por \textbf{C2}, el mismo par de estados $\langle p, p' \rangle$. Por hipótesis de inducción, las subderivaciones que cubren esas dos ocurrencias producen subcadenas $u^1$ y $u^2$ tales que $\transit(p, u^1) = p'$ y $\transit(p, u^2) = p'$. La igualdad $u^1 = u^2$ no se sigue del autómata por sí solo, ya que en general puede haber múltiples cadenas distintas que conecten dos estados. La igualdad sí se sigue de la semántica de $G$: una cláusula heredada en $G'$ conserva las variables originales en las mismas posiciones, y al aplicar la cláusula original $c$ en $G$, la variable $X$ se asigna a un único rango y por tanto a una única subcadena. Las dos subderivaciones de $G'$ producen, vistas como derivaciones de $G$, la misma subcadena para $X$, lo cual cierra la consistencia.

Aplicada la cláusula original $c$ con la asignación de rangos así determinada, $A$ deriva en $G$ la tupla $(u_1, \dots, u_a)$. La condición (ii) de la Afirmación~\ref{afirm:induccion} se obtiene por concatenación de las propiedades de pares de las subderivaciones, que se ensamblan correctamente gracias a la condición \textbf{C1}.

La Afirmación~\ref{afirm:induccion} queda demostrada. Aplicada al símbolo $S[\langle \initial, q \rangle]$ derivando $w$, se obtiene $w \in \Lang(G)$, lo cual completa la dirección $\Lang(G') \subseteq \Lang(G) \cap \Lang(A)$.

\medskip
\noindent\textbf{Inclusión $\Lang(G) \cap \Lang(A) \subseteq \Lang(G')$.}\\
Sea $w \in \Lang(G) \cap \Lang(A)$. Por la pertenencia a $\Lang(G)$ existe una derivación $D$ en $G$ de $S(\langle 0, \lvert w \rvert \rangle)$ hacia $\eps$. Por la pertenencia a $\Lang(A)$ se cumple $\transit(\initial, w) = q$ con $q \in \final$. Por la Observación~\ref{obs:recorrido}, el recorrido de $A$ sobre $w$ está dado por la sucesión única $s_0, s_1, \dots, s_{\lvert w \rvert}$ con $s_0 = \initial$, $s_{\lvert w \rvert} = q$, y $s_k = \transit(\initial, w[1 \ldots k])$.

Se construye una derivación de $w$ en $G'$ replicando la estructura de $D$ y indexando cada predicado con los pares de estados correspondientes. La asignación es la siguiente: si en $D$ una subderivación parte de un predicado $A(\langle i_1, j_1 \rangle, \dots, \langle i_a, j_a \rangle)$, en $G'$ se indexa $A$ con el vector $\vec{q} = (\langle s_{i_1}, s_{j_1} \rangle, \dots, \langle s_{i_a}, s_{j_a} \rangle)$. Análogamente, si $D$ asigna a una variable $X$ el rango $\langle a, b \rangle$, se le asocia el par $\langle s_a, s_b \rangle$.

Esta asignación satisface la condición \textbf{C1} porque, para cada argumento $\alpha_{j,i}$ de la cláusula aplicada en $D$, los terminales y las posiciones límite de las variables coinciden con segmentos del recorrido de $A$ sobre $w$, y las transiciones inducidas por los terminales son exactamente las transiciones del autómata sobre $w$. Satisface la condición \textbf{C2} porque, en $D$, dos ocurrencias de una misma variable se asignan al mismo rango por la semántica de \PRCG, y por tanto al mismo par de estados por la unicidad del recorrido.

Por la Definición~\ref{def:gprime}, las cláusulas heredadas resultantes están en $\Pset'$. La derivación en $G'$ que se obtiene reemplazando cada predicado de $D$ por su versión indexada conduce de $S'(\langle 0, \lvert w \rvert \rangle)$ hacia $\eps$, comenzando con la cláusula inicial $S'(X) \to S[\langle \initial, q \rangle](X)$, de donde $w \in \Lang(G')$.

La doble inclusión completa la demostración.
\end{proof}

% =============================================================
\subsection{Tamaño de la gramática producto y un ejemplo}\label{sec:tamano}

\subsubsection{Tamaño de \texorpdfstring{$G'$}{Gprime} y complejidad del algoritmo}\label{sec:complejidad}

\paragraph{Cantidad de no terminales.} El conjunto de no terminales $\Nset'$ tiene cardinalidad acotada por $\lvert \Nset \rvert \cdot \lvert \statesT \rvert^{2 a_{\max}} + 1$, donde $a_{\max}$ es la aridad máxima de los predicados de $G$. La cota se obtiene observando que cada predicado $A$ de aridad $a$ produce $\lvert \statesT \rvert^{2 a}$ versiones indexadas, una por cada vector de pares de estados.

\paragraph{Cantidad de cláusulas.} El conjunto de cláusulas $\Pset'$ tiene cardinalidad acotada por
\[
\lvert \Pset' \rvert \;\leq\; \sum_{c \in \Pset} \lvert \statesT \rvert^{2 n_c},
\]
donde $n_c$ es la cantidad de variables distintas que aparecen en la cláusula $c$. La cota se obtiene observando que cada asignación $\Phi$ de pares de estados a las $n_c$ variables distintas de $c$ produce a lo sumo una cláusula heredada, y que el espacio de asignaciones $\Phi$ tiene cardinalidad $\lvert \statesT \rvert^{2 n_c}$. Las condiciones \textbf{C1} y \textbf{C2} filtran ese espacio: \textbf{C2} se cumple por construcción al modelar $\Phi$ sobre variables distintas, y \textbf{C1} elimina las asignaciones inconsistentes con las transiciones del autómata. La cota es siempre válida como techo y se alcanza en gramáticas con muchas variables y autómatas en los que todas las asignaciones son consistentes.

\paragraph{Complejidad temporal del algoritmo.} El cuerpo principal del algoritmo recorre cada cláusula $c \in \Pset$ y, para cada $c$, enumera el espacio de asignaciones $\Phi$. Cada asignación se verifica frente a \textbf{C1} en tiempo proporcional al largo total de los argumentos de $c$, denotado $\ell_c$. La cantidad de cláusulas heredadas producidas para $c$ está acotada por $\lvert \statesT \rvert^{2 K_c}$, donde $K_c$ es la suma de las aridades de los predicados de $c$, ya que para cada $\Phi$ que pasa \textbf{C1} se enumeran las combinaciones de pares para los argumentos de cada predicado. El costo total del algoritmo es por tanto
\[
O\!\left( \sum_{c \in \Pset} \lvert \statesT \rvert^{2 n_c} \cdot \ell_c \cdot \lvert \statesT \rvert^{2 K_c} \right).
\]
Acotando uniformemente sobre las cláusulas, sean $n_{\max} = \max_c n_c$, $K_{\max} = \max_c K_c$ y $\ell_{\max} = \max_c \ell_c$. La complejidad queda
\[
O\!\left( \lvert \Pset \rvert \cdot \ell_{\max} \cdot \lvert \statesT \rvert^{2(n_{\max} + K_{\max})} \right).
\]
La dependencia en $\lvert \statesT \rvert$ es polinomial de grado $2(n_{\max} + K_{\max})$. La dependencia en $\lvert \Pset \rvert$ y en $\ell_{\max}$ es lineal.

\paragraph{Polinomialidad respecto al autómata.} Para una \PRCG\ $G$ fijada, el algoritmo opera en tiempo polinomial respecto al tamaño del autómata. El grado del polinomio depende de la aridad y de la cantidad de variables de las cláusulas, parámetros propios de $G$. Esta dependencia es la misma que la de la construcción de Bar-Hillel para gramáticas libres del contexto, donde el grado es dos. La aparición del factor $K_{\max}$ corresponde a la generalización a predicados de aridad mayor, y la aparición del factor $n_{\max}$ corresponde a la enumeración de asignaciones $\Phi$ que la condición \textbf{C2} requiere.

\paragraph{Comportamiento cuando \texorpdfstring{$G$}{G} es entrada del problema.} Si la \PRCG\ no se fija y se considera entrada del problema junto con el autómata, la dependencia es exponencial en la aridad y en la cantidad de variables. La cota inferior de Bertsch y Nederhof~\cite{REF}, que demuestra la $\mathbf{NP}$-completitud del problema del vacío de la intersección cuando el autómata reconoce un lenguaje regular finito, indica que tal dependencia exponencial no se puede evitar en general, salvo si $\mathbf{P} = \mathbf{NP}$. La construcción del Teorema~\ref{thm:correctness} produce $G'$ explícitamente, lo cual implica que el tamaño de $G'$ está acotado por el tiempo de cómputo del algoritmo y por tanto no puede ser menor que las cotas anteriores.

\subsubsection{Ejemplo trabajado}

Se retoma $G_0$ y $A_0$ del ejemplo~\ref{sec:idea}. La construcción produce los siguientes símbolos no terminales: el nuevo distinguido $S'$; las nueve versiones $S[\langle q_i, q_j \rangle]$ del predicado de aridad uno $S$ original, una por cada par $\langle q_i, q_j \rangle \in \statesT \times \statesT$; y las nueve versiones $\mathrm{Par}[\langle q_i, q_j \rangle]$ del predicado de aridad uno $\mathrm{Par}$ original. En total, $1 + 9 + 9 = 19$ símbolos no terminales.

Las cláusulas iniciales son las que conectan $S'$ con las versiones de $S$ que terminan en un estado final. Como $\final = \{q_0\}$, solo se conserva
\begin{Verbatim}[xleftmargin=2em]
  S'(X) -> S[<q0, q0>](X).
\end{Verbatim}

Las cláusulas heredadas a partir de $c_1 : S(X) \to \mathrm{Par}(X) \; \mathrm{Par}(X)$ se construyen así. La cláusula tiene una variable $X$ con tres ocurrencias: una en la cabeza y dos en el cuerpo, una por cada predicado $\mathrm{Par}$. La condición \textbf{C2} fuerza que las tres ocurrencias reciban el mismo par. Sea $\Phi(X) = \langle p, p' \rangle$ ese par único. La condición \textbf{C1} aplicada al argumento de la cabeza, $\alpha = X$, da inmediatamente que el par del argumento de la cabeza coincide con $\langle p, p' \rangle$. La cláusula instanciada es
\[
S[\langle p, p' \rangle](X) \to \mathrm{Par}[\langle p, p' \rangle](X) \; \mathrm{Par}[\langle p, p' \rangle](X).
\]

\begin{tcolorbox}[colback=blue!5!white, colframe=blue!50!black, boxrule=0.5pt, arc=2pt]
\textbf{Punto clave.} La condición \textbf{C2} impone que los dos predicados $\mathrm{Par}$ del cuerpo lleven el mismo par de estados que la cabeza. Sin \textbf{C2}, los dos $\mathrm{Par}$ podrían recibir pares distintos, lo que permitiría cláusulas como $S[\langle q_0, q_0 \rangle](X) \to \mathrm{Par}[\langle q_0, q_1 \rangle](X) \; \mathrm{Par}[\langle q_1, q_0 \rangle](X)$, donde la misma $X$ tendría que recorrer dos pares distintos del autómata simultáneamente, lo cual genera cadenas que no están en $\Lang(G_0) \cap \Lang(A_0)$. La verificación operacional confirma este punto en la sección~\ref{sec:literatura}.
\end{tcolorbox}

Como $\Phi(X)$ recorre los nueve pares posibles, se obtienen nueve cláusulas heredadas a partir de $c_1$. De ellas, solamente $S[\langle q_0, q_0 \rangle](X) \to \mathrm{Par}[\langle q_0, q_0 \rangle](X) \; \mathrm{Par}[\langle q_0, q_0 \rangle](X)$ es alcanzable desde la cláusula inicial.

Las cláusulas heredadas a partir de $c_2 : \mathrm{Par}(aaX) \to \mathrm{Par}(X)$ se construyen del modo siguiente. La cabeza tiene un único argumento $aaX$. Sea $\Phi(X) = \langle p, p' \rangle$. La condición \textbf{C1} aplicada al argumento $\alpha = aaX$, descompuesto como $\beta_1 = a, \beta_2 = a, \beta_3 = X$, exige una sucesión de estados $r_0 = p_0, r_1, r_2, r_3 = p'_0$ tal que $\trans(r_0, a) = r_1$, $\trans(r_1, a) = r_2$, y $\langle r_2, r_3 \rangle = \langle p, p' \rangle$. De aquí $r_2 = p$ y $r_3 = p'$, así que $\trans(\trans(r_0, a), a) = p$. Es decir, $r_0$ es el predecesor de $p$ por dos transiciones de $a$. En el autómata $A_0$, dos transiciones de $a$ corresponden al avance cíclico $q_0 \to q_2$, $q_1 \to q_0$, $q_2 \to q_1$, así que el predecesor doble se calcula:
\[
\text{pred}^2(q_0) = q_1, \qquad \text{pred}^2(q_1) = q_2, \qquad \text{pred}^2(q_2) = q_0.
\]
El esquema de cláusula heredada es
\begin{Verbatim}[xleftmargin=2em]
  Par[<pred^2(p), p'>](aaX) -> Par[<p, p'>](X).
\end{Verbatim}
Como $\Phi(X)$ recorre los nueve pares posibles, se obtienen nueve cláusulas heredadas a partir de $c_2$.

La cláusula terminal $c_3 : \mathrm{Par}(\eps) \to \eps$ se trata como sigue. El argumento $\eps$ obliga a que el par asignado satisfaga $p_0 = p'_0$, ya que la cadena vacía no produce transición. Hay tres cláusulas heredadas, una por cada estado de $\statesT$:
\begin{Verbatim}[xleftmargin=2em]
  Par[<q0, q0>](eps) -> eps,
  Par[<q1, q1>](eps) -> eps,
  Par[<q2, q2>](eps) -> eps.
\end{Verbatim}

En total, $G'$ tiene $19$ símbolos no terminales y, antes de eliminar los símbolos no productivos, $1 + 9 + 9 + 3 = 22$ cláusulas: una inicial, nueve heredadas a partir de $c_1$, nueve a partir de $c_2$ y tres terminales a partir de $c_3$.

Después de eliminar los símbolos que no derivan en ninguna cadena, queda una \PRCG\ reducida que reconoce únicamente las cadenas $a^{2n}$ con $2n$ múltiplo de $3$, esto es, cadenas $a^{6m}$, que es $\Lang(G_0) \cap \Lang(A_0)$.

El ejemplo deja ver dos cosas. Primero, la cantidad de cláusulas y de símbolos crece según las cotas de la sección~\ref{sec:tamano}. Segundo, la condición \textbf{C2} acopla las ocurrencias de $X$ entre la cabeza y los dos predicados $\mathrm{Par}$ del cuerpo, eliminando combinaciones de pares que no respetan la igualdad semántica de la variable. Sin ese acoplamiento, la gramática producto admitiría cláusulas heredadas adicionales que generan cadenas fuera de $\Lang(G_0) \cap \Lang(A_0)$, como se documenta en la sección~\ref{sec:literatura}.

% =============================================================
\subsection{Posición respecto a la literatura y verificación operacional}\label{sec:literatura}

Esta sección sitúa la construcción del Teorema~\ref{thm:correctness} frente a las construcciones publicadas que tratan problemas relacionados, y describe la verificación operacional realizada para corroborar la corrección.

\subsubsection{Relación con construcciones publicadas}

La idea de indexar predicados con pares de estados para construir el cierre bajo intersección con regulares aparece en la literatura para subclases más restringidas. Bar-Hillel, Perles y Shamir~\cite{REF} la introducen para gramáticas libres del contexto. Seki, Matsumura, Fujii y Kasami~\cite{REF} la extienden a gramáticas multi-contextuales y a su variante paralela mediante tuplas de pares de estados, en el Teorema~3.9(3) de su trabajo de 1991. Bertsch y Nederhof~\cite{REF} la formulan en el caso particular de las \sRCG\ y obtienen un algoritmo que opera en tiempo polinomial respecto al tamaño combinado de la gramática y del autómata.

La construcción del presente capítulo se diferencia de las anteriores en dos aspectos. El primero es el formalismo objetivo: las \PRCG\ admiten cláusulas en las que una misma variable aparece en varias posiciones, característica que es inalcanzable en \sRCG\ por la restricción de linealidad. El segundo es la aparición de la condición \textbf{C2} como pieza estructural de la construcción. En \sRCG\ la condición \textbf{C2} no es necesaria porque las cláusulas son lineales y ninguna variable aparece más de una vez. En el caso PMCFG, donde la copia se manifiesta en la función de composición de las reglas y no en la repetición de variables, la prueba de Seki~\cite{REF} aborda el problema con una técnica distinta cuyo desarrollo detallado para el caso copia no figura explícitamente en el texto de 1991.

En las búsquedas realizadas en la literatura revisada por pares y en la literatura gris consultada no aparece una construcción explícita análoga a la del Teorema~\ref{thm:correctness} para la clase de las \PRCG\ top-down non-erasing. Boullier~\cite{REF} establece el cierre de las RCG bajo intersección mediante una construcción declarativa que añade una cláusula de conjunción y no construye una gramática producto con pares de estados. Boullier y Sagot~\cite{REF} adaptan el parsing de RCG a entradas representadas por DAGs acíclicos, pero no producen una gramática producto, y su garantía de tiempo polinomial se restringe a las \sRCG.

La construcción del presente capítulo cubre una clase estrictamente más expresiva que las \sRCG. La \sRCG\ se define como gramática lineal, non-erasing y no-combinatorial en el sentido de la Definición~8 de Boullier~\cite{REF}; la clase de las \PRCG\ top-down non-erasing es estrictamente más amplia porque admite no-linealidad y combinatorialidad. Boullier~\cite{REF} demuestra además que la propiedad de linealidad no se preserva por equivalencia: existen lenguajes generados por \PRCG\ no-lineales que no admiten una \PRCG\ lineal equivalente. Por tanto, la construcción del presente capítulo cubre lenguajes que están fuera del alcance de cualquier algoritmo de intersección sobre \sRCG, incluido el de Bertsch y Nederhof~\cite{REF}.

\subsubsection{El papel de la semántica de PRCG en la condición \texorpdfstring{\textbf{C2}}{C2}}

Una pregunta natural es si la condición \textbf{C2}, que iguala los pares de estados asociados a dos ocurrencias de la misma variable, basta para que la gramática producto reconozca exactamente la intersección. La pregunta es legítima porque, en un \AFD\ genérico, dos cadenas distintas pueden conectar el mismo par de estados, lo que en principio podría permitir que la gramática producto generara una cadena en la que dos ocurrencias de una variable correspondieran a subcadenas distintas pero con el mismo recorrido en $A$.

La razón por la que \textbf{C2} es suficiente reside en la semántica de las \PRCG. En una derivación, una variable se asigna a un único rango sobre la cadena de entrada, y por tanto a una única subcadena. La gramática producto $G'$ hereda esa semántica porque las cláusulas heredadas conservan literalmente las variables en las posiciones de la cláusula original; al aplicar una cláusula heredada en $G'$, la misma variable, repetida en posiciones distintas, recibe un único rango. La condición \textbf{C2} sirve para que esos pares de estados sean coherentes; la igualdad de las subcadenas es una consecuencia de la semántica de la \PRCG, no del autómata.

Este punto se contrasta con el caso de la PMCFG. En PMCFG la copia ocurre a través de la función de composición que duplica una variable como argumento de salida, y la ``asignación de rango único por variable'' no se aplica del mismo modo. Por esa razón, una construcción de tipo Seki para PMCFG con copias requiere un análisis de coherencia distinto al de \textbf{C2}, y allí radica la dificultad técnica que la literatura no resuelve de forma explícita en el texto original.

\subsubsection{Necesidad de \texorpdfstring{\textbf{C2}}{C2}: un contraejemplo mínimo}\label{sec:contraejemplo-c2}

Se exhibe a continuación un contraejemplo explícito que justifica la necesidad de la condición \textbf{C2}. La existencia del contraejemplo demuestra que la omisión de \textbf{C2}, manteniendo \textbf{C1}, produce una gramática producto cuyo lenguaje contiene cadenas que no pertenecen a $\Lang(G) \cap \Lang(A)$. La construcción es, por tanto, incorrecta sin \textbf{C2}.

\paragraph{La gramática $G_c$.} Sea $G_c$ la \PRCG\ con predicados $S$ y $P$ de aridad uno, sobre el alfabeto $\Tnt = \{a\}$, y las cláusulas siguientes:
\[
\begin{aligned}
c_1 &: \quad S(X) \to P(X) \; P(X) \\
c_2 &: \quad P(aX) \to P(X) \\
c_3 &: \quad P(\eps) \to \eps
\end{aligned}
\]
La gramática $G_c$ es top-down non-erasing en el sentido de la sección~\ref{sec:hipotesis}: en $c_1$ la variable $X$ aparece en la cabeza y en los dos predicados del cuerpo; en $c_2$ la variable $X$ aparece en la cabeza y en el cuerpo; en $c_3$ no hay variables. Las cláusulas $c_2$ y $c_3$ definen $P$ como el conjunto de cadenas sobre $\{a\}$. La cláusula $c_1$ exige que la cadena entera sea aceptada simultáneamente por dos invocaciones de $P$ sobre el mismo argumento $X$. Por tanto $\Lang(G_c) = \{a^n : n \geq 0\}$.

\paragraph{El autómata $A_c$.} Sea $A_c$ el \AFD\ con dos estados $\{q_0, q_1\}$, alfabeto $\{a\}$, estado inicial $q_0$, estado final $q_0$, y transiciones $\trans(q_0, a) = q_1$ y $\trans(q_1, a) = q_0$. El autómata $A_c$ reconoce las cadenas de longitud par sobre $\{a\}$, esto es, $\Lang(A_c) = \{a^{2k} : k \geq 0\}$.

\paragraph{La intersección.} La intersección $\Lang(G_c) \cap \Lang(A_c)$ es el conjunto de cadenas de longitud par sobre $\{a\}$:
\[
\Lang(G_c) \cap \Lang(A_c) = \{a^{2k} : k \geq 0\}.
\]

\paragraph{Construcción correcta con \textbf{C1} y \textbf{C2}.} La construcción de la Definición~\ref{def:gprime} produce una gramática $G'$ de $11$ cláusulas. Las cláusulas heredadas de $c_1$ son cuatro:
\[
\begin{aligned}
S[\langle q_0, q_0 \rangle](X) &\to P[\langle q_0, q_0 \rangle](X) \; P[\langle q_0, q_0 \rangle](X) \\
S[\langle q_0, q_1 \rangle](X) &\to P[\langle q_0, q_1 \rangle](X) \; P[\langle q_0, q_1 \rangle](X) \\
S[\langle q_1, q_0 \rangle](X) &\to P[\langle q_1, q_0 \rangle](X) \; P[\langle q_1, q_0 \rangle](X) \\
S[\langle q_1, q_1 \rangle](X) &\to P[\langle q_1, q_1 \rangle](X) \; P[\langle q_1, q_1 \rangle](X)
\end{aligned}
\]
En cada una de las cuatro cláusulas, las dos invocaciones de $P$ reciben el \emph{mismo} par de estados. Esa coincidencia es la materialización de \textbf{C2}: la variable $X$, al recibir un único par bajo $\Phi$, transmite el mismo par a las dos ocurrencias. La gramática $G'$ reconoce exactamente $\{a^{2k} : k \geq 0\}$.

\paragraph{Construcción incorrecta sin \textbf{C2}.} Si la construcción se modifica para asignar un par independiente a cada \emph{ocurrencia} sintáctica de $X$ en $c_1$ y se conserva la condición \textbf{C1}, las cláusulas heredadas de $c_1$ son $64$ en lugar de cuatro. Cada vector de cuatro pares de estados, uno por cada ocurrencia sintáctica de $X$ en la cláusula, produce una cláusula heredada distinta. Entre ellas aparece, por ejemplo,
\[
S[\langle q_0, q_0 \rangle](X) \to P[\langle q_0, q_1 \rangle](X) \; P[\langle q_1, q_0 \rangle](X).
\]
Esta cláusula afirma que la cadena $X$ recorre simultáneamente el autómata $A_c$ desde $q_0$ hasta $q_1$ y desde $q_1$ hasta $q_0$. Ninguna cadena fija puede cumplir ambas exigencias a la vez, pero la cláusula heredada las acepta como compatibles porque \textbf{C1} solo verifica la consistencia local de cada argumento, no la coherencia global entre las ocurrencias. La gramática producto resultante, denotada $G'_{\overline{c2}}$, contiene $83$ cláusulas. La cadena $a$ pertenece a $\Lang(G'_{\overline{c2}})$: aplicando la cláusula inicial $S'(X) \to S[\langle q_0, q_0 \rangle](X)$ con $X = \langle 0, 1 \rangle$, luego la cláusula heredada anterior, y observando que la subcadena $a$ va de $q_0$ a $q_1$ en $A_c$ y también de $q_1$ a $q_0$ en $A_c$, las dos invocaciones de $P$ se derivan independientemente sobre la misma cadena $a$. El resultado es que $a \in \Lang(G'_{\overline{c2}})$.

\paragraph{Demostración de la sobre-generación.} La cadena $a$ pertenece a $\Lang(G_c)$ porque $G_c$ genera todas las cadenas sobre $\{a\}$, pero no pertenece a $\Lang(A_c)$ porque su longitud es impar. Por tanto $a \notin \Lang(G_c) \cap \Lang(A_c)$. Como $a \in \Lang(G'_{\overline{c2}})$, la gramática producto sin \textbf{C2} no satisface la igualdad de lenguajes que el Teorema~\ref{thm:correctness} establece. El mismo argumento se aplica a todas las cadenas $a^{2k+1}$ con $k \geq 0$: la versión sin \textbf{C2} reconoce $\Lang(G_c) = \Sigma^{*}$, no la intersección $\{a^{2k} : k \geq 0\}$.

\paragraph{Conclusión.} La condición \textbf{C2} no es una restricción opcional. Su omisión, manteniendo \textbf{C1}, produce sobre-generación demostrable sobre una \PRCG\ top-down non-erasing de tres cláusulas y un \AFD\ de dos estados. La presencia de \textbf{C2} en la construcción de la Definición~\ref{def:gprime} corresponde al ajuste mínimo necesario para que la construcción Bar-Hillel-style se extienda correctamente desde \sRCG\ hasta la clase top-down non-erasing tratada en este capítulo.

\subsubsection{Verificación operacional}

Se implementó la construcción de la Definición~\ref{def:gprime} en una rutina de referencia, junto con un intérprete de \PRCG\ capaz de decidir la pertenencia de una cadena al lenguaje generado. Para examinar la suficiencia de la condición \textbf{C2}, se implementó adicionalmente una variante de la rutina que aplica únicamente la condición \textbf{C1} y omite \textbf{C2}, asignando un par de estados a cada \emph{ocurrencia} sintáctica de cada variable de manera independiente, en lugar de un par único por variable. La existencia de discrepancias entre las dos versiones permite contrastar empíricamente la necesidad de \textbf{C2} mediante un \emph{differential test}: si la versión sin \textbf{C2} produce el mismo lenguaje que la versión con \textbf{C2} sobre todos los casos, entonces \textbf{C2} no aporta restricciones adicionales sobre esos casos; si produce un lenguaje estrictamente mayor, entonces \textbf{C2} es responsable de eliminar la sobre-generación.

Se diseñó un conjunto de casos de prueba que incluye \PRCG\ con distintas formas de no-linealidad: no-linealidad bottom-up con argumentos distintos en la cabeza, no-linealidad top-down con la misma variable en varios predicados del cuerpo, y combinaciones. Se incluyeron también gramáticas lineales según Boullier~\cite{REF} que generan el lenguaje de copia $\{ww : w \in \{a, b\}^{*}\}$, los lenguajes $a^n b^n$ y $a^n b^n c^n$, y casos límite ($\Lang(A) = \emptyset$, $\Lang(A) = \{\eps\}$, gramáticas que reconocen únicamente $\eps$). Para cada combinación de $G$ y $A$, se enumeraron exhaustivamente las cadenas de longitud acotada y se comparó la pertenencia a $\Lang(G')$ con la pertenencia simultánea a $\Lang(G)$ y a $\Lang(A)$.

\paragraph{Resultados con la construcción correcta.}

En todos los casos probados, la versión que aplica \textbf{C1} y \textbf{C2} satisface $\Lang(G') = \Lang(G) \cap \Lang(A)$ sobre las cadenas enumeradas. Esto incluye intersecciones con autómatas adversariales en los que múltiples cadenas distintas conectan los mismos pares de estados, donde un fallo en la coherencia entre ocurrencias produciría falsos positivos detectables por la enumeración. Esos casos no produjeron ningún falso positivo.

\paragraph{Resultados sin \textbf{C2}.}

La versión que omite \textbf{C2} sobre-genera en varios casos. Las observaciones más relevantes son las siguientes.

\begin{itemize}
\item Sobre la \PRCG\ $G_0$ del ejemplo motivador, $S(X) \to \mathrm{Par}(X) \; \mathrm{Par}(X)$, intersectada con el autómata $A_0$ de longitud múltiplo de tres, la versión con \textbf{C2} produce $22$ cláusulas y reconoce exactamente $\{a^{6m}\}$. La versión sin \textbf{C2} produce $814$ cláusulas y reconoce además las cadenas $a^2$ y $a^4$, que pertenecen a $\Lang(G_0)$ pero no a $\Lang(A_0)$.
\item Sobre la cláusula $S(X) \to P(X) \; Q(X)$ con $P, Q$ predicados que reconocen $a^{*}$, intersectada con el autómata de longitud múltiplo de tres, la versión sin \textbf{C2} sobre-genera con las cadenas $a, a^2, a^4, a^5$.
\item Sobre la cláusula $S(X) \to P(X) \; Q(X) \; R(X)$ con tres predicados, intersectada con autómatas de longitud múltiplo de dos y de tres, la versión sin \textbf{C2} sobre-genera con cadenas explícitas en cada caso.
\item Sobre una \PRCG\ \emph{lineal} en el sentido de Boullier que reconoce $a^{*}$, mediante la cláusula $S(aX) \to S(X)$ junto a $S(\eps) \to \eps$, intersectada con el autómata de longitud múltiplo de tres, la versión con \textbf{C2} produce $13$ cláusulas y reconoce exactamente $\{a^{3k}\}$, mientras que la versión sin \textbf{C2} produce $85$ cláusulas y sobre-genera con $a, a^2, a^4, a^5, a^7$.
\end{itemize}

\paragraph{Lectura de los resultados.}

La observación sobre la gramática lineal $S(aX) \to S(X)$ tiene un alcance que conviene comentar. La cláusula es lineal en el sentido de la Definición~8 de Boullier~\cite{REF}: la variable $X$ aparece exactamente una vez en la cabeza y exactamente una vez en el cuerpo. Sin embargo, la versión sin \textbf{C2} de la construcción produce sobre-generación incluso en este caso, porque al asignar pares por ocurrencia las dos apariciones de $X$ pueden recibir pares incompatibles, lo que rompe la coherencia entre la cabeza y el cuerpo. Por tanto, \textbf{C2} no es una restricción específica para \PRCG\ no-lineales, sino una pieza estructural necesaria de la construcción Bar-Hillel-style para cualquier \PRCG\ top-down non-erasing donde una variable aparece en más de una posición de la cláusula. Lo que distingue las \PRCG\ top-down non-erasing tratadas en este capítulo de las \sRCG\ no es la existencia de \textbf{C2}, sino la cantidad y diversidad de configuraciones donde \textbf{C2} actúa de manera no trivial: en \sRCG\ una variable puede aparecer a lo sumo una vez en cabeza y una vez en cuerpo, y \textbf{C2} acopla únicamente esas dos ocurrencias; en una \PRCG\ top-down non-erasing una variable puede aparecer cualquier número de veces, y \textbf{C2} acopla todas las ocurrencias simultáneamente.

La verificación operacional no constituye prueba formal ---no puede sustituir la demostración del Teorema~\ref{thm:correctness}--- pero proporciona dos pieces de evidencia empírica relevantes: la versión con \textbf{C1} y \textbf{C2} no presenta sobre-generación en ningún caso probado, y la versión sin \textbf{C2} produce sobre-generación documentada en múltiples casos no triviales, incluyendo gramáticas lineales en el sentido de Boullier.

% =============================================================
\subsection{Discusión y transición al capítulo siguiente}\label{sec:transicion}

El Teorema~\ref{thm:correctness} cierra la pregunta de la construcción de la intersección de una \PRCG\ top-down non-erasing con un lenguaje regular mediante un objeto explícito que se puede inspeccionar. El resultado no debe confundirse con un algoritmo de decisión para el vacío de $\Lang(G) \cap \Lang(A)$. De acuerdo con Bertsch y Nederhof~\cite{REF}, ese problema es indecidible para \PRCG\ generales y $R$ regular cualquiera, y se vuelve $\mathbf{NP}$-completo si se restringe $R$ a lenguajes finitos.

En el capítulo siguiente se aborda la pregunta planteada por los tutores: qué propiedades estructurales de las \PRCG\ son responsables de esa indecidibilidad y qué subclase se obtiene al imponerlas. La construcción presentada en este capítulo será la base sobre la cual se mide el efecto de imponer linealidad o no-combinatorialidad sobre la hipótesis top-down non-erasing ya adoptada: bajo cada restricción adicional, la gramática producto $G'$ adquiere propiedades que pueden hacer decidible el vacío y, en algunos casos, incluso polinomial.

En particular, cuando $G$ cumple los tres criterios de simplicidad de las \sRCG\ ---linealidad, non-erasure y no-combinatorialidad---, la cantidad de configuraciones donde la condición \textbf{C2} actúa de manera no trivial es limitada: una variable aparece a lo sumo una vez en la cabeza y una vez en el cuerpo, y \textbf{C2} acopla únicamente esas dos ocurrencias. La construcción de la sección~\ref{sec:construccion}, restringida a ese caso, coincide esencialmente con la construcción polinomial que Bertsch y Nederhof~\cite{REF} dan para \sRCG, lo que permite decidir el vacío de $\Lang(G')$ en tiempo polinomial respecto al tamaño de $G'$. La observación delimita el rol de la no-linealidad como propiedad responsable de la indecidibilidad y abre el camino al análisis del capítulo siguiente, en el que se examinan las restricciones que recuperan la decidibilidad sin colapsar al caso simple.

\end{document}
